% Retained preamble for this source, sliced MECHANICALLY from its own main TeX.
% Source file : full-original-native.tex (lines 1-77 of the file)
% The block below is verbatim author text; the wrapper class, document metadata and
% \begin{document} are omitted because the assembler supplies them.
% lines taken: 1, 2, 3, 4, 5, 6, 7, 8, 9, 10, ... (77 lines total)
%\documentclass[JEP,XML,SOM,Unicode,NoFloatCountersInSection]{cedram}
\datereceived{2019-04-04}
\dateaccepted{2019-11-18}
\dateepreuves{2019-12-03}

\usepackage{mathrsfs}
\let\mathcal\mathscr
\let\epsilon\varepsilon
\let\sep\mid
\newcommand{\lrangle}{\langle\cdot\,,\cdot\rangle}

\newenvironment{enumeratea}
{\bgroup\def\theenumi{\alph{enumi}}\begin{enumerate}}
{\end{enumerate}\egroup}

\newcommand{\spfrac}[2]{\sfrac{#1}{(#2)}}
\newcommand{\psfrac}[2]{\sfrac{(#1)}{#2}}
\newcommand{\pspfrac}[2]{\sfrac{(#1)}{(#2)}}
\newcommand{\Psfrac}[2]{(\sfrac{#1}{#2})}
\newcommand{\Ppsfrac}[2]{(\sfrac{(#1)}{#2})}

\newcommand{\Changel}{\mathcode`l="7160}
\newcommand{\Changelback}{\mathcode`l="716C}
\newcommand{\NoChangel}[1]{%
\expandafter\let\csname old\string#1\endcsname=#1
\let#1=\relax
\newcommand{#1}{\mathcode`l="716C\csname old\string#1\endcsname\mathcode`l="7160 }%
}
\def\mto{\mathchoice{\longmapsto}{\mapsto}{\mapsto}{\mapsto}}
\newcommand{\cpr}{c\mkern1mu'}
\newcommand{\bcpr}{\bar c\mkern1mu'}
\newcommand{\tcpr}{\tilde c\mkern2mu'}
\newcommand{\sts}{\mathrm{ss}}
\newcommand{\uu}{\mathrm{uu}}

\NoChangel{\log}
\NoChangel{\ln}
\NoChangel{\lim}
\NoChangel{\limsup}
\NoChangel{\liminf}
\NoChangel{\varlimsup}
\NoChangel{\varliminf}
\NoChangel{\varinjlim}
\NoChangel{\varprojlim}

\usepackage{mathtools}
\usepackage{color}
\usepackage{tikz}
\usetikzlibrary{patterns}

\newtheorem{thm}{Theorem}[section]
\newtheorem{thmx}{Theorem}
\renewcommand{\thethmx}{\Alph{thmx}}
\newtheorem{corx}[thmx]{Corollary}
\newtheorem{lem}[thm]{Lemma}
\newtheorem{prop}[thm]{Proposition}
\newtheorem{cor}[thm]{Corollary}

\theoremstyle{definition}
\newtheorem{defn}[thm]{Definition}

\numberwithin{equation}{section}

\newcommand\R{\ifmmode{\mathbb R}\else{$\mathbb R$}\fi}

\DeclareMathOperator{\CAT}{CAT}
\DeclareMathOperator{\diam}{diam}
\DeclareMathOperator{\Int}{Int}
\DeclareMathOperator{\Proj}{Proj}
\DeclareMathOperator{\Susp}{Susp}

\newcommand\DDD{\mathcal D}
\newcommand\KKK{\mathcal K}
\newcommand\BBB{\mathcal B}
\newcommand\RRR{\mathcal R}
\newcommand\SSS{\mathcal S}

